\documentclass{article}

\usepackage{geometry}
\geometry{a4paper}

\usepackage[round,authoryear]{natbib}
\bibliographystyle{plainnat}
\usepackage[hidelinks]{hyperref}
\usepackage{authblk}
\usepackage{tabularray}
\usepackage{graphicx}
\usepackage{amssymb}

% Make affiliations run inline
\renewcommand\Affilfont{\small}
\renewcommand\Authfont{\normalsize}

\makeatletter
\renewcommand\AB@affilsepx{;\hspace{0.6em}\protect\Affilfont}
\makeatother

\renewcommand{\thetable}{S\arabic{table}}
\renewcommand{\thefigure}{S\arabic{figure}}
\renewcommand{\thesubsection}{\arabic{subsection}}

\begin{document}

\title{Population-scale Ancestral Recombination Graphs with tskit 1.0\\
\large Supplementary Information}

\input{authors.tex}

\date{}
\maketitle

\subsection*{Published tools using tskit}
We performed a search for published software using tskit.
To be included in this list, the software must be described in a publication
as a named tool intended for reuse and must use tskit.
This list is intentionally conservative, and does not include
(for example) publications depending on bespoke analysis
pipelines using tskit.
For each tool in Table S1, we list: its name (with a clickable
hyperlink to its source code repository);
the main implementation
language(s) of the tool based on GitHub summaries;
the DOI (with clickable hyperlink) of the publication in which this
tool was first described (or, in the case of SLiM, the first
publication in which tskit-based features were described);
and the year of publication and number of citations to this
publication according to data from OpenAlex
(\url{https://openalex.org}) retrieved on 2026-01-23.

\input{tools_table.tex}

\subsection*{Tskit functionality and benchmarks}
To give a sense of the breadth of functionality provided by tskit,
Table S2 lists properties and operations it supports, grouped
into nine broad categories.
There is no direct competitor to tskit, so a head-to-head feature
comparison is not possible. Instead, we contrast it with three
Python libraries from adjacent niches:
ARGneedle-lib~\citep{zhang2023biobank}, which infers and analyses
ancestral recombination graphs from biobank-scale data;
the BTE interface (\url{https://github.com/jmcbroome/BTE})
for matUtils~\citep{mcbroome2021daily},
a toolkit for manipulating and analysing mutation-annotated trees
of SARS-CoV-2; and DendroPy~\citep{sukumaran2010dendropy}, a
long-established general-purpose phylogenetic scripting library.
For each operation we mark whether the comparison library offers a
fully comparable feature (\checkmark), a partial or restricted
analogue ($\circ$), or none (blank). Marks reflect library state
at the time of writing.

Tskit is the result of more than a decade of sustained development,
and its distinctive operations have been introduced in a
series of publications. The succinct tree
sequence data format and the sequential tree-construction algorithm
were first introduced as part of the msprime coalescent
simulator~\citep{kelleher2016efficient}.
The key algorithmic result is
that, given the edge table of an ARG, the cost of iterating over
the marginal trees along the genome scales with the number of trees
and the edge differences between neighbouring trees, rather than
with the sample size.
The practical consequence was demonstrated in a benchmark
where using the succinct tree sequence format and sequential
iteration was more than a million times faster than reading the
same trees individually from Newick using the standard libraries
available at the time, for a simulation of 100,000 samples.
The data model was subsequently generalised to enable
forwards-time simulation~\citep{kelleher2018efficient}. This
publication
also introduced the columnar tables API and was the first paper in
which the tskit library itself appeared as a distinct component.
That paper gave a rigorous time- and space-complexity analysis of
the ``simplify'' operation, which restricts an ARG to the
history ancestral to a specified set of samples. ARG simplification
is a fundamental operation with many applications~\citep{wong2024general},
as  well as being the primitive operation making forwards-time simulation of ARGs
tractable~\citep{haller2018tree}.

Following this foundational work, tskit
accumulated a sequence of new operations, each first described in its own
publication. Motivated by needing to summarise biobank-scale inferred ARGs
with uncalibrated node times, \citet{kelleher2019inferring} introduced the
the genealogical nearest neighbours (GNN) statistic.
\citet{ralph2020efficiently} then introduced a general API for
population-genetic statistics built around a duality theorem
relating site-based and branch-based statistics.
The benchmarks in that paper show
speedups of up to five orders of magnitude against
genotype-matrix implementations at a fraction of the memory footprint.
Individual and pedigree abstractions were added to the tskit data
model by \citet{andersontrocme2023genes}, extending what a node or
individual in an ARG can represent.
\citet{tsambos2023link} introduced the ``link\_ancestors"
operation, a new primitive for extracting local-ancestry tracts
from an ARG.

[TODO this paragraph needs heavy revision, it's quite weak.]
A further class of operation became possible once the core
traversal orders were exposed to Numba~\citep{lam2015numba},
allowing Python-level user code to JIT-compile ARG traversals at
close to C speed. This enabled downstream packages
like tstrait~\citep{tagami2024tstrait}
and tsbrowse~\citep{karthikeyan2025tsbrowse} to perform large-scale
computations using only Python code.
\citet{fritze2026forest} added the
``extend\_haplotypes'' function to tskit.
\citet{lehmann2026on}
introduced the operation of computing genetic relatedness matrices
(GRMs) directly from an ARG. \citet{lee2026genetic} apply this
last operation to genomic prediction and provide the one recent
head-to-head benchmark of tskit against another library,
comparing the ARG-based GRM with the corresponding implementation
in ARGneedle-lib. Tskit~1.0 is the point at which this
decade-long accumulation of operations is brought together and
stabilised under a single, guaranteed API. The many blank cells
in Table S2 are not an artefact of a narrow choice of comparison
libraries: most of the operations listed above were defined by the
publications that introduced them and simply have no analogue
elsewhere, which is also why a direct head-to-head benchmark suite
against a competing library is not possible.

\input{functionality_table.tex}

\bibliography{paper}

\end{document}
